\chapter{The Framework of Nearly Everything}
\label{chap:framework-of-everything}

We are now in possession of all the pieces of our cosmological puzzle: what fundamentally exists (ontology), what is typical (measure) and
what then inevitably emerges (the structure).

\section{The Ontology: What Fundamentally Exists}

\subsection*{Undefined Unknown}

We do not assume finiteness.

If reality were fundamentally finite, one could immediately ask: finite by how much? Why n bits rather than n+1? What principle fixed that particular bound?

Any fixed finite limit appears contingent and therefore calls for a deeper explanatory layer.

For this reason, we do not posit a fundamentally finite substrate.


We also do not assume information as primitive.

Information presupposes distinguishability. At its simplest, this reduces to binary distinction: a bit, a difference between alternatives.

But declaring reality to be fundamentally informational or digital is itself a substantive ontological commitment.

A discrete ontology is conceivable, but so is a continuous one. Nothing prior licenses choosing one baseline over the other.

Therefore, we do not presuppose that reality is fundamentally finite, discrete, informational, mathematical, or governed by prior, external laws.

We denote this unrestricted, pre-interpretive totality ($\mathbf{U}$). We do not claim to know what it is.

Only that no specific formatting principle is granted ontological privilege a priori.


\subsection*{The Observer Constraint}

Despite radical agnosticism about the substrate, one empirical fact remains unavoidable: observer exists. Conscious, structured  experience exists.

Based on all observational evidence, we as human beings possess a strictly finite informational structure.

As derived in the Chapter~\nameref{chap:interpretation-problem} information is not an inherent property of objective reality, but an emergent artifact of observation.

The ``bit'' is not the fundamental building block of the cosmos; it is emergent unit of distinction employed by a finite interpreter.

As derived in the Chapter~\nameref{chap:humans-as-axiomatic-systems} the universe is static.
Under this view, no metaphysical signals transmit information from outside to inside observer.

It then follows that every possible experience accessible to an observer must already be encoded within the total observer structure.


\subsection*{The Guitar String Analogy}

To explain how finite discreteness can emerge from infinite an intuitive analog is a guitar string.

The string itself is perfectly continuous. It isn't made of discrete "tone particles."

Discreteness resonant modes emerge from the boundary conditions: when you pluck it, it can only vibrate at discrete,
quantized frequencies (the fundamental note, the octave, the third harmonic, etc.), 

This is because the string is clamped at both ends! The boundaries force a continuous medium to manifest discrete states.

Likewise, a finite observer interacting with an unrestricted reality may induce discrete informational structure without discreteness being ontologically primitive.

Under algorithmic or Solomonoff-style measures, low-complexity, highly compressible structures dominate typicality.

Finite observers would therefore be expected to find themselves embedded within stable, structured, compressible equivalence classes.

This motivates the minimal ontological assumption:

\[
\text{Undefined} \rightarrow  \text{Finite Observer} 
\]

The step remains unexplained. That is an explicit boundary of the theory.

We do not know what the unrestricted undefined totality is.

We assume only that finite observers exist within it.

Because observers are finite, any observer-accessible description can be encoded informationally without assuming that the underlying ontology itself is informational.

Information is not what reality is made of; information is what finite observers perceive when interacting with an otherwise unrestricted reality.


\section{Measure: What is Typical}


\subsection*{The Initial Assumption: Configuration Space}

For formal descriptive purposes, we introduce a finite, static configuration space:
\[
\mathcal{C} = \{0,1\}^n
\]
where each $c \in \mathcal{C}$ is a complete, static snapshot of bits.
Bits play no fundamental role; the binary system is chosen just for the sake of convenience.

Each configuration $c$ admits multiple potential descriptions $d \in \mathcal{D}(c)$.
We quantify the total complexity $C_{\text{total}}(d)$ of a description by splitting it into its dual spectral and geometric representations:
\[
C_{\text{total}}(d) = C_s(d) + C_g(d)
\]

The spectral complexity $C_s(d)$ measures the algorithmic bit-cost required to track the wave-like modes:
\[
C_s(d) = \sum_i \left[ \text{Cost}(A_i) + \text{Cost}(\phi_i) + \left(\frac{\omega_i}{\Delta \omega}\right) \right]
\]
where the sum runs over active spectral modes, $\text{Cost}(A_i)$ represents the bit-depth of the amplitude envelope, $\text{Cost}(\phi_i)$ is the fixed-width register encoding the periodic phase interval $[0, 2\pi]$, and $\omega_i / \Delta \omega$ represents the linear resource allocation required to track the mode frequency relative to the minimum resolution threshold $\Delta\omega$.

The geometric complexity term $C_g(d)$ accounts for localized spatial structures, boundaries, and metric fields.



\subsection*{The Equation of Joint Compressibility}

Because the relationship between these two is a many-to-many mapping, any given observer wavefunction can be interpreted through nearly infinitely many different geometric spacetime metrics, and conversely, any spacetime manifold can be decomposed into nearly infinitely many different wavefunctions.
This creates an unimaginably massive configuration space. 

In information theory, when a system must be simultaneously validated across two dual representations,
its joint probability is the product of its individual probabilities.

Therefore, the absolute measure of any emergent physical state ($s$) within the totality is given by the \textbf{Joint Compressibility Equation}:

\begin{equation}
P_{\text{joint}}(s) = P_{\text{spectral}}(s) \times P_{\text{geometric}}(s) = 2^{-C_s(s)} \times 2^{-C_g(s)} = 2^{-\left(C_s(s) + C_g(s)\right)}
\end{equation}

This deceptively simple formula yields a profound variational principle for a unified framework.
Nature minimizes the total description length.
The cosmos is governed by a singular, overarching informational Lagrangian:

\begin{equation}
\label{eq:lagrangian}
\mathcal{L}_{\text{info}} = C_s(s) + C_g(s) \longrightarrow \text{Minimum}
\end{equation}

Neither geometry nor waves are primary; they are equal, co-dependent projections of a single underlying informational matrix.
Our observed reality is the sweet spot of this joint optimization.

Spacetime is smooth and the quantum world is lawful because this specific configuration represents the absolute shortest combined code nature can write.



\subsection*{Induced Measure and Observer Functionals}

The raw algorithmic weight of a static configuration $c$ is accumulated across all its valid descriptions:
\[
W(c) = \sum_{d \in \mathcal{D}(c)} 2^{-C_{\text{total}}(d)}
\]
This defines our raw baseline probability distribution over the state space:
\[
P(c) = \frac{W(c)}{\sum_{c'\in\mathcal{C}} W(c')}
\]

Because $C_s(d)$ maps complexity linearly to frequency, the calculation $2^{-C_{\text{total}}(d)}$ inside the summation guarantees that high-frequency states suffer genuine exponential probability suppression, mirroring the Boltzmann distribution ($P \propto e^{-\beta E}$) found in statistical mechanics.

An observer is modeled as a grading semantic functional:
\[
\mu_O : \mathcal{C} \to [0,1]
\]
which measures how strongly a given bit configuration describes observer-relative structure.
This functional induces an observer-relative equivalence relation $c_1 \sim_O c_2$ whenever $c_1$ and $c_2$ describe
the same relational observer structure to within the tolerances encoded by $\mu_O$. 

The observer-conditioned probability then becomes:
\[
P(c \mid O) = \frac{\mu_O(c)W(c)}{\sum_{c'} \mu_O(c')W(c')}
\]


\subsection*{The Mathematical Unification}

To transition from the probability of a single static configuration $P(c \mid O)$ to the probability of an entire experienced history,
we must track how an observer traverses the configuration space.
Let a path $\gamma \in \Gamma_O$ represent a sequence of configurations ordered by our internal temporal relation $\prec_O$.
Because the configurations along the path are independent slices of the static totality,
the total probability of the path is the product of its individual constituent measures, meaning their algorithmic costs sum additively.

By changing our informational units from bits to natural units (where $2^{-C} = e^{-C \ln 2}$), we can express the global probability of an entire observed history $\gamma$:
\[
\mathbb{P}(\gamma \mid O) = \frac{1}{Z_O} \exp\left(-\lambda\mathcal{C}_O[\gamma]\right), \qquad \gamma \in \Gamma_O
\]
where $\mathcal{C}_O[\gamma] = \int_{\gamma} C_{\text{total}}(c)\, dt$ represents the accumulated informational cost along the path, and $\lambda$ acts as an algorithmic inverse temperature (scaling the bit-to-nat conversion and the observer's filtering acuity).
This states that observed physical laws are simply the large-deviation minimizers of the total informational cost function $\mathcal{C}_O[\gamma]$ over the set of all semantic paths.



\section{Emergence: The Structure}

\subsection*{The Emergence of Spacetime and Laws}

Our core hypothesis is that configurations capable of supporting observers strongly favor smooth, low-frequency spectral content.
Sharp, chaotic particle trajectories require massive high-frequency modes, which exponentially spike the linear description length $C_s(d)$ and are instantly suppressed by the joint measure. 

As a result, observer-compatible configurations are overwhelmingly dominated by those exhibiting inertia, smooth curved paths, and law-like behavior.
Spacetime geometry emerges because it makes macroscopic informational boundaries computationally cheap to describe ($C_g$). 

It emerges entirely as an internal ordering relation $\prec_O$ on configurations, where $c_1 \prec_O c_2$ if $c_2$ supports the observer functional $\mu_O$.
The experienced flow of time corresponds to traversal along the most compressible, optimized paths through this timeless superposition.



\subsection*{Trinity and the Minimal Resolution}

We have three emergent layers of reality:

\[
D-\psi-G
\]

\begin{itemize}
    \item $\mathbf{D}$ \textbf{(Data/Description):} The raw, naked, discrete digital—manifesting physically as localized particles or pixels. 
    \item $\boldsymbol{\psi}$ \textbf{(Spectral/Quantum):} The smooth wave-codec that compresses the relational possibility space, minimizing spectral cost ($C_s$).
    \item $\mathbf{G}$ \textbf{(Geometry/Gravity):} The continuous metric manifold that compresses the realization space at scale, minimizing geometric boundaries ($C_g$).
\end{itemize}

Discrete reality (particles) are localized, high-density ``knots'' where the smooth continuous wave-codec ($\psi$) meets the finite digital register of the observer.
The global physics remains governed entirely by our dual informational Lagrangian:

\begin{equation}
\mathcal{L}_{\text{info}} = C_s(s) + C_g(s) \longrightarrow \text{Minimum}
\end{equation}

The discrete data ($D$) is the emergent \emph{output} of this joint minimization when parsed by a bounded interpreter.
Spacetime is smooth ($G$), fields are continuous ($\psi$), yet our observations are invariably quantized into discrete particless ($D$).




\section{The Equations}


\begin{tcolorbox}[title=\iameverything]

\[
\mathbb{P}(\gamma \mid O) = \frac{1}{Z_O} \exp\left(-\lambda\mathcal{C}_O[\gamma]\right), \qquad \gamma \in \Gamma_O
\]

  
\[
D-\psi-G
\]


\end{tcolorbox}
